\chapter{경보 신호}
% full version: chapters/14_pain.tex

\pencilfig{14}{\pencilart{14}}{통증은 온도계가 아니라 화재경보기다.}

다른 모든 감각은 세상이 어떤지를 알린다. 통증은 그렇지 않다. 통증은 무언가가 몸을 위협한다고 알리며, 그 크기는 자극보다 기계 자신이 정하는 부분이 크다. 통증은 온도계가 아니라 화재경보기다.

통증 감지기는 강한 것에만 반응한다. 40여 도가 넘는 열, 강한 압력, 부식성 물질이다. 통증 감지기는 촉각 감지기보다 훨씬 덜 순응한다. 위협이 있는 동안 경보가 멈추면 안 되기 때문이다.

\section*{전선 두 개}

손가락을 어딘가에 부딪쳤던 때를 떠올려 보자. 먼저 날카롭고 정확한 통증이 온다. 어디서, 언제인지를 알린다. 잠시 뒤 둔하고 퍼지는 긴 통증이 온다. 얼마나 나쁜지를 알린다. 이 둘은 서로 다른 전선이다. 하나는 빨라서 신호가 수십 밀리초 만에 도착한다. 다른 하나는 느려서 신호가 1초 정도 걸린다.

\section*{관문}

부딪친 곳을 우리는 왜 문지를까? 척수에서 통증 신호는 “관문”을 지난다. 촉각은 관문을 반쯤 닫는 억제성 뉴런을 깨운다. 그래서 부딪친 곳을 문지르면 실제로 도움이 된다. 다만 중간 정도의 통증에만 그렇고, 강한 통증은 어쨌든 통과한다. 강한 통증이 이때 억제성 뉴런을 스스로 끈다는 것은 원래의 관문 이론의 일부이며, 아직 직접 확인되지 않았다.

\pencilfig{126}{%
% gate: the pain wire drives the relay neuron; touch wakes an inhibitory neuron that closes the gate
\draw[pen=pcAmber] (0,1.35) -- (3.05,1.0); \draw[arr=pcAmber] (2.8,1.03) -- (3.08,1.0);
\node[lbl, text=pcAmber!70!black, anchor=east] at (-0.05,1.35) {통증};
\draw[pen=pcBlue] (0,-0.15) -- (1.75,-0.15); \draw[arr=pcBlue] (1.5,-0.15) -- (1.78,-0.15);
\node[lbl, text=pcBlue, anchor=east] at (-0.05,-0.15) {촉각};
\draw[dotpen=pcRed, wash=pcRed] (2.05,-0.15) circle (0.26);
\draw[pcRed, line width=0.7pt, -{Bar[width=6pt]}] (2.25,0.05) -- (3.2,0.66);
\node[lbl, text=pcRed, anchor=north] at (2.05,-0.45) {억제};
\draw[dotpen=pcGray, wash=pcGray] (3.4,0.9) circle (0.34);
\node[lbl, anchor=south] at (3.4,1.3) {관문};
\draw[pen] (3.75,0.9) -- (5.6,0.9); \draw[arr] (5.35,0.9) -- (5.65,0.9);
\node[lbl, anchor=west] at (5.7,0.9) {뇌로};
}{척수의 관문. 통증 신호는 관문 뉴런을 거쳐 뇌로 가고, 촉각은 관문을 반쯤 닫는 억제성 뉴런을 깨운다. 그래서 부딪친 곳을 문지르면 도움이 되지만, 중간 정도의 통증에만 그렇다.}

\section*{위에서 돌리는 음량 노브}

뇌는 통증을 스스로 키우거나 줄일 수 있다. 위쪽의 뇌에서 관문으로 라디오 채널이 내려온다. \nt{SERO}, \nt{NORA}, 그리고 뇌가 스스로 분비하는, 진통제와 닮은 특별한 물질이다. 위험한 상황에서는 통증이 거의 사라질 수 있다. 위험에서 도망치는 사람에게는 절뚝거릴 여유가 없다. 통증이 줄어들 것이라는 기대도 통증을 누그러뜨리며, 일부는 같은 물질을 통해서다.

\section*{경보가 들러붙을 때}

통증이 되풀이되면 민감도가 올라간다. 같은 충격이 더 아프게 느껴진다. 이것도 학습이다. 상처가 아무는 동안 스스로를 조여 가는 루프다. 그러나 강한 루프에는 위험이 있다. 루프가 “들러붙어”, 상처가 다 나은 뒤에도 계속 울릴 수 있다. 2장에서 본, 스스로를 유지하는 뉴런 무리와 같다.

\section*{교사이자 인터럽트}

통증은 교사이기도 하다. 도파민의 “예상보다 좋은가, 나쁜가”에 대해 통증은 음의 보상으로 작용한다. 또 통증은 노르아드레날린의 분출처럼 하던 일을 중단시킨다. 모든 것을 내려놓고 위험부터 처리하라. 그리고 가장 빠른 반응, 곧 뜨거운 것에서 손을 떼는 반응은 앞 장에서 본 바로 그 맞선 근육 한 쌍을 통해 척수 안에서 곧바로 닫힌다. 여러분이 통증을 느끼기도 전에 말이다.

\fullversion{14}
